\documentclass[10pt,a4paper]{amsart}
\usepackage{amsmath,amssymb,amsthm,graphicx}
% Retained preamble for the Stacks Project chapter snapshot in this directory.
% The chapter's own preamble.tex is a copy of the shared Stacks Project preamble; the block
% below is sliced MECHANICALLY, line for line, from it. Nothing was added, removed or
% redefined; the wrapper class, the xr-hyper/hyperref wiring and the externaldocument list
% are omitted because the assembler supplies the document class.
%
% lines taken from preamble.tex: 3, 5, 9, 11, 12, 156, 157, 158, ... (51 lines total)
%
% package coverage restored from the chapter's own preamble (2026-10-13): xy, xy-2cell

\usepackage[all]{xy}
\xyoption{2cell}
\UseAllTwocells

\newenvironment{reference}{\comment}{\endcomment}
\newenvironment{slogan}{\comment}{\endcomment}
\newenvironment{history}{\comment}{\endcomment}
\theoremstyle{plain}
\newtheorem{theorem}[subsection]{Theorem}
\newtheorem{proposition}[subsection]{Proposition}
\newtheorem{lemma}[subsection]{Lemma}
\theoremstyle{definition}
\newtheorem{definition}[subsection]{Definition}
\newtheorem{example}[subsection]{Example}
\newtheorem{exercise}[subsection]{Exercise}
\newtheorem{situation}[subsection]{Situation}
\theoremstyle{remark}
\newtheorem{remark}[subsection]{Remark}
\newtheorem{remarks}[subsection]{Remarks}
\numberwithin{equation}{subsection}
\def\lim{\mathop{\mathrm{lim}}\nolimits}
\def\colim{\mathop{\mathrm{colim}}\nolimits}
\def\Spec{\mathop{\mathrm{Spec}}}
\def\Hom{\mathop{\mathrm{Hom}}\nolimits}
\def\Ext{\mathop{\mathrm{Ext}}\nolimits}
\def\SheafHom{\mathop{\mathcal{H}\!\mathit{om}}\nolimits}
\def\SheafExt{\mathop{\mathcal{E}\!\mathit{xt}}\nolimits}
\def\Sch{\mathit{Sch}}
\def\Mor{\mathop{\mathrm{Mor}}\nolimits}
\def\Ob{\mathop{\mathrm{Ob}}\nolimits}
\def\Sh{\mathop{\mathit{Sh}}\nolimits}
\def\NL{\mathop{N\!L}\nolimits}
\def\CH{\mathop{\mathrm{CH}}\nolimits}
\def\proetale{{pro\text{-}\acute{e}tale}}
\def\etale{{\acute{e}tale}}
\def\QCoh{\mathit{QCoh}}
\def\Ker{\mathop{\mathrm{Ker}}}
\def\Im{\mathop{\mathrm{Im}}}
\def\Coker{\mathop{\mathrm{Coker}}}
\def\Coim{\mathop{\mathrm{Coim}}}
\DeclareMathSymbol{\boxtimes}{\mathbin}{AMSa}{"02}
\def\QCohstack{\mathcal{QC}\!\mathit{oh}}
\def\Cohstack{\mathcal{C}\!\mathit{oh}}
\def\Spacesstack{\mathcal{S}\!\mathit{paces}}
\def\Quotfunctor{\mathrm{Quot}}
\def\Hilbfunctor{\mathrm{Hilb}}
\def\Curvesstack{\mathcal{C}\!\mathit{urves}}
\def\Polarizedstack{\mathcal{P}\!\mathit{olarized}}
\def\Complexesstack{\mathcal{C}\!\mathit{omplexes}}
\def\Pic{\mathop{\mathrm{Pic}}\nolimits}
\def\Picardstack{\mathcal{P}\!\mathit{ic}}
\def\Picardfunctor{\mathrm{Pic}}
\def\Deformationcategory{\mathcal{D}\!\mathit{ef}}
\usepackage{cleveref}
\crefname{originallemma}{original lemma}{original lemmas}
\Crefname{originallemma}{Original lemma}{Original lemmas}
\makeatletter
\newlabel{algebra-lemma-Artin-Rees}{{00IN}{source}{}{sourcepointer-1}{}}
\newlabel{algebra-lemma-Artin-Rees@cref}{{[originallemma][0][]00IN}{[1][0][]source}{}{}{}}
\makeatother
\crefname{originallemma}{original lemma}{original lemmas}
\Crefname{originallemma}{Original lemma}{Original lemmas}
\makeatletter
\newlabel{example-key}{{07YI}{source}{}{sourcepointer-2}{}}
\newlabel{example-key@cref}{{[originallemma][0][]07YI}{[1][0][]source}{}{}{}}
\makeatother
\crefname{originallemma}{original lemma}{original lemmas}
\Crefname{originallemma}{Original lemma}{Original lemmas}
\makeatletter
\newlabel{formal-defos-lemma-essential-surjection}{{06H0}{source}{}{sourcepointer-3}{}}
\newlabel{formal-defos-lemma-essential-surjection@cref}{{[originallemma][0][]06H0}{[1][0][]source}{}{}{}}
\makeatother
\crefname{originallemma}{original lemma}{original lemmas}
\Crefname{originallemma}{Original lemma}{Original lemmas}
\makeatletter
\newlabel{formal-defos-lemma-versal-criterion}{{06IU}{source}{}{sourcepointer-4}{}}
\newlabel{formal-defos-lemma-versal-criterion@cref}{{[originallemma][0][]06IU}{[1][0][]source}{}{}{}}
\makeatother
\crefname{originallemma}{original lemma}{original lemmas}
\Crefname{originallemma}{Original lemma}{Original lemmas}
\makeatletter
\newlabel{lemma-compute-ext-into-field}{{07YK}{source}{}{sourcepointer-5}{}}
\newlabel{lemma-compute-ext-into-field@cref}{{[originallemma][0][]07YK}{[1][0][]source}{}{}{}}
\makeatother
\crefname{originallemma}{original lemma}{original lemmas}
\Crefname{originallemma}{Original lemma}{Original lemmas}
\makeatletter
\newlabel{lemma-construct-essential-surjection}{{07YL}{source}{}{sourcepointer-6}{}}
\newlabel{lemma-construct-essential-surjection@cref}{{[originallemma][0][]07YL}{[1][0][]source}{}{}{}}
\makeatother
\crefname{originallemma}{original lemma}{original lemmas}
\Crefname{originallemma}{Original lemma}{Original lemmas}
\makeatletter
\newlabel{remark-canonical-element}{{07YC}{source}{}{sourcepointer-7}{}}
\newlabel{remark-canonical-element@cref}{{[originallemma][0][]07YC}{[1][0][]source}{}{}{}}
\makeatother
\title{Versality characterized by the cotangent complex comparison}
\author{The Stacks Project Authors}
\date{stacks-project:artin.tex:local-snapshot-20261012; GNU FDL 1.2 or later}
\begin{document}
\maketitle
\noindent\textit{Editorial note.} Statement span is lines 4234-4265 of sources/stacks-artin/pinned-original-source.tex; proof span is lines 4267-4357, the complete author proof (89 body lines).\par\medskip
\noindent\textit{Editorial note.} Navigation stubs: 7 cited-but-not-included labels are supplied as non-load-bearing pointers with their pinned tags where the snapshot tag index has them.\par\medskip
\noindent\textit{Editorial note.} extra\_packages is empty on purpose: the retained preamble already loads xy with options and a second declaration clashes in the assembler wrapper.\par\medskip
\noindent\textit{Editorial note.} Text is retained verbatim, including the author's own typos and placeholders.\par\medskip
\noindent\textit{Editorial note.} Queue index 6 of queue\_artin.json (pb=4066, stmt 4234-4265, proof 4267-4357).\par\medskip
\noindent\textit{Editorial note.} stacks-artin wave ar01 (2026-10-14): fresh adjudication accept per lane-s3/artin\_judgments.md; pool per Chief 10:5x direct dispatch (idle-first over S12's 01:5x note, SMC precedent); artin.tex 99 rows/67 with proof judged 46 accept / 21 reject / 5 AGP-overlap avoidance rows; patch-sd-envs contract spans scanned.\par\medskip
\noindent\textit{Allowed standard prerequisites (editorial).}\begin{itemize}
\item Explicit prerequisite boundary. The cited inputs are prerequisites; the argument in the retained span is not replaced by any cited result.
\end{itemize}\medskip
\ifdefined\hypertarget\hypertarget{sourcepointer-1}{}\fi
\noindent\textit{Source reference.} Original lemma 00IN: sources/stacks-artin/pinned-original-source.tex, source lines 4357--4357. Cited by the retained statement or proof but not part of the collected spans; supplied only as a navigation pointer. The argument does not depend on its content.\par\medskip
\ifdefined\hypertarget\hypertarget{sourcepointer-2}{}\fi
\noindent\textit{Source reference.} Original lemma 07YI: sources/stacks-artin/pinned-original-source.tex, source lines 4068--4091. Cited by the retained statement or proof but not part of the collected spans; supplied only as a navigation pointer. The argument does not depend on its content.\par\medskip
\ifdefined\hypertarget\hypertarget{sourcepointer-3}{}\fi
\noindent\textit{Source reference.} Original lemma 06H0: sources/stacks-artin/pinned-original-source.tex, source lines 4357--4357. Cited by the retained statement or proof but not part of the collected spans; supplied only as a navigation pointer. The argument does not depend on its content.\par\medskip
\ifdefined\hypertarget\hypertarget{sourcepointer-4}{}\fi
\noindent\textit{Source reference.} Original lemma 06IU: sources/stacks-artin/pinned-original-source.tex, source lines 4357--4357. Cited by the retained statement or proof but not part of the collected spans; supplied only as a navigation pointer. The argument does not depend on its content.\par\medskip
\ifdefined\hypertarget\hypertarget{sourcepointer-5}{}\fi
\noindent\textit{Source reference.} Original lemma 07YK: sources/stacks-artin/pinned-original-source.tex, source lines 4150--4153. Cited by the retained statement or proof but not part of the collected spans; supplied only as a navigation pointer. The argument does not depend on its content.\par\medskip
\ifdefined\hypertarget\hypertarget{sourcepointer-6}{}\fi
\noindent\textit{Source reference.} Original lemma 07YL: sources/stacks-artin/pinned-original-source.tex, source lines 4161--4173. Cited by the retained statement or proof but not part of the collected spans; supplied only as a navigation pointer. The argument does not depend on its content.\par\medskip
\ifdefined\hypertarget\hypertarget{sourcepointer-7}{}\fi
\noindent\textit{Source reference.} Original lemma 07YC: sources/stacks-artin/pinned-original-source.tex, source lines 3806--3838. Cited by the retained statement or proof but not part of the collected spans; supplied only as a navigation pointer. The argument does not depend on its content.\par\medskip
\begin{lemma}
\label{lemma-characterize-versal}
Let $S$ be a locally Noetherian scheme. Let $\mathcal{X}$ be a category
fibred in groupoids over $(\Sch/S)_{fppf}$ satisfying (RS*).
Let $U = \Spec(A)$ be an
affine scheme of finite type over $S$ which maps into an affine open
$\Spec(\Lambda)$. Let $x$ be an object of $\mathcal{X}$ over $U$.
Let $\xi : E \to \NL_{A/\Lambda}$ be a morphism of $D^{-}(A)$. Assume
\begin{enumerate}
\item[(i)] for every deformation situation $(x, A' \to A)$ we have:
$x$ lifts to $\Spec(A')$ if and only if
$E \to \NL_{A/\Lambda} \to \NL_{A/A'}$ is zero, and
\item[(ii)] there is an isomorphism of functors
$T_x(-) \to \Ext^0_A(E, -)$
such that $E \to \NL_{A/\Lambda} \to \Omega^1_{A/\Lambda}$
corresponds to the canonical element (see
Remark \ref{remark-canonical-element}).
\end{enumerate}
Let $u_0 \in U$ be a finite type point with residue field
$k = \kappa(u_0)$. Consider the following statements
\begin{enumerate}
\item $x$ is versal at $u_0$, and
\item $\xi : E \to \NL_{A/\Lambda}$ induces a surjection
$H^{-1}(E \otimes_A^{\mathbf{L}} k) \to
H^{-1}(\NL_{A/\Lambda} \otimes_A^{\mathbf{L}} k)$
and an injection
$H^0(E \otimes_A^{\mathbf{L}} k) \to
H^0(\NL_{A/\Lambda} \otimes_A^{\mathbf{L}} k)$.
\end{enumerate}
Then we always have (2) $\Rightarrow$ (1) and we have (1) $\Rightarrow$ (2)
if $u_0$ is a closed point.
\end{lemma}

\begin{proof}
Let $\mathfrak p = \Ker(A \to k)$ be the prime corresponding to $u_0$.

\medskip\noindent
Assume that $x$ versal at $u_0$ and that $u_0$ is a closed point of $U$.
If $H^{-1}(\xi \otimes_A^{\mathbf{L}} k)$ is not surjective, then
let $A' \to A$ be an extension with kernel $I$ as in
Lemma \ref{lemma-construct-essential-surjection}.
Because $u_0$ is a closed point, we see that $I$ is a finite $A$-module,
hence that $A'$ is a finite type $\Lambda$-algebra (this fails if
$u_0$ is not closed). In particular $A'$ is Noetherian.
By property (c) for $A'$ and (i) for $\xi$ we see that $x$ lifts to
an object $x'$ over $A'$.
Let $\mathfrak p' \subset A'$ be kernel of the surjective map to $k$.
By Artin-Rees (Algebra, Lemma \ref{algebra-lemma-Artin-Rees})
there exists an $n > 1$ such that $(\mathfrak p')^n \cap I = 0$.
Then we see that
$$
B' = A'/(\mathfrak p')^n \longrightarrow A/\mathfrak p^n = B
$$
is a small, essential extension of local Artinian rings, see
Formal Deformation Theory, Lemma
\ref{formal-defos-lemma-essential-surjection}.
On the other hand, as $x$ is versal at $u_0$ and as $x'|_{\Spec(B')}$
is a lift of $x|_{\Spec(B)}$, there exists an integer
$m \geq n$ and a map $q : A/\mathfrak p^m \to B'$
such that the composition
$A/\mathfrak p^m \to B' \to B$ is the quotient map.
Since the maximal ideal of $B'$ has $n$th power equal to zero, this
$q$ factors through $B$ which contradicts the fact that $B' \to B$ is an
essential surjection. This contradiction shows that
$H^{-1}(\xi \otimes_A^{\mathbf{L}} k)$
is surjective.

\medskip\noindent
Assume that $x$ versal at $u_0$. By Lemma \ref{lemma-compute-ext-into-field}
the map $H^0(\xi \otimes_A^{\mathbf{L}} k)$ is dual to the map
$\Ext^0_A(\NL_{A/\Lambda}, k) \to \text{Ext}^0_A(E, k)$. Note that
$$
\Ext^0_A(\NL_{A/\Lambda}, k) = \text{Der}_\Lambda(A, k)
\quad\text{and}\quad
T_x(k) = \Ext^0_A(E, k)
$$
Condition (ii) assures us the map
$\Ext^0_A(\NL_{A/\Lambda}, k) \to \text{Ext}^0_A(E, k)$
sends a tangent vector $\theta$ to $U$ at $u_0$ to the corresponding
infinitesimal deformation of $x_0$, see Remark \ref{remark-canonical-element}.
Hence if $x$ is versal, then this map is surjective, see
Formal Deformation Theory, Lemma \ref{formal-defos-lemma-versal-criterion}.
Hence $H^0(\xi \otimes_A^{\mathbf{L}} k)$ is injective.
This finishes the proof of (1) $\Rightarrow$ (2) in case $u_0$ is a
closed point.

\medskip\noindent
For the rest of the proof assume $H^{-1}(E \otimes_A^\mathbf{L} k) \to
H^{-1}(\NL_{A/\Lambda} \otimes_A^\mathbf{L} k)$
is surjective and
$H^0(E \otimes_A^\mathbf{L} k) \to
H^0(\NL_{A/\Lambda} \otimes_A^\mathbf{L} k)$
injective. Set $R = A_\mathfrak p^\wedge$ and let $\eta$ be the
formal object over $R$ associated to $x|_{\Spec(R)}$.
The map $d\underline{\eta}$ on tangent spaces is surjective
because it is identified with the dual of the injective map
$H^0(E \otimes_A^{\mathbf{L}} k) \to
H^0(\NL_{A/\Lambda} \otimes_A^{\mathbf{L}} k)$
(see previous paragraph). According to
Formal Deformation Theory, Lemma \ref{formal-defos-lemma-versal-criterion}
it suffices to prove the following:
Let $C' \to C$ be a small extension of finite type Artinian local
$\Lambda$-algebras with residue field $k$. Let $R \to C$ be a
$\Lambda$-algebra map compatible with identifications of residue fields.
Let $y = x|_{\Spec(C)}$ and let $y'$ be a lift of $y$ to $C'$.
To show: we can lift the $\Lambda$-algebra map $R \to C$ to $R \to C'$.

\medskip\noindent
Observe that it suffices to lift the $\Lambda$-algebra map $A \to C$.
Let $I = \Ker(C' \to C)$. Note that $I$ is a $1$-dimensional $k$-vector
space. The obstruction $ob$ to lifting $A \to C$ is an element of
$\Ext^1_A(\NL_{A/\Lambda}, I)$, see Example \ref{example-key}.
By Lemma \ref{lemma-compute-ext-into-field} and our assumption the map
$\xi$ induces an injection
$$
\Ext^1_A(\NL_{A/\Lambda}, I)
\longrightarrow
\Ext^1_A(E, I)
$$
By the construction of $ob$ and (i) the image of $ob$ in $\Ext^1_A(E, I)$
is the obstruction to lifting $x$ to $A \times_C C'$. By (RS*) the fact that
$y/C$ lifts to $y'/C'$ implies that $x$ lifts to $A \times_C C'$. Hence
$ob = 0$ and we are done.
\end{proof}

\end{document}
